% ─────────────────────────────────────────────────────────────────────────────
\section{Caching Strategy}
% ─────────────────────────────────────────────────────────────────────────────

\subsection{Motivation}

Two separate caching concerns arise in \catena.

\textbf{Tail-convergent caching.}  The recurrence
\[
    h_n = g(n)\,h_{n-1} + h_{n-2}, \quad k_n = g(n)\,k_{n-1} + k_{n-2}
\]
has optimal substructure: $h_n$ depends on $h_{n-1}$ and $h_{n-2}$.
Without memoisation a naive recursive implementation recomputes all earlier
tail convergents each time a later one is requested.  With memoisation each
value is computed once, reducing the total cost from $O(2^n)$ to $O(n)$.
All \texttt{SimpleContinuedFraction} instances carry a dedicated
\texttt{\_tail\_cache} for this purpose.

\textbf{Generator caching.}  When the partial-quotient function $g: \N \to \Zp$
is expensive to evaluate (e.g.\ a number-theoretic function or a transcendental
constant computed digit by digit), re-evaluating $g(n)$ on every recurrence step
would dominate the runtime.  \texttt{CachedGenerator} wraps an arbitrary callable
and memoises each distinct index in a \texttt{BaseCache}.

\subsection{The \texttt{BaseCache} Primitive}

\catena's caching layer is built around a single primitive defined in
\texttt{catena.cache}: \texttt{BaseCache}.  It is a self-computing,
append-only dictionary (\texttt{UserDict} subclass) that unifies the three
separate classes that existed in earlier versions (\texttt{Cache},
\texttt{OrdinalCache}, and \texttt{CacheHandler}).

\begin{description}[style=nextline]
  \item[Immutability after write.]
    Once a key is stored, any subsequent attempt to overwrite it raises
    \texttt{KeyError}.  This is correct for convergents and generator values
    alike, because both are deterministic functions of their index for a fixed SCF.

  \item[Self-computing on miss.]
    \texttt{BaseCache} is constructed with a bound callable \texttt{func}.
    Accessing a missing key via \texttt{cache[key]} calls \texttt{func(key)},
    stores the result, and returns it — no external decorator or wrapper is
    needed.  Explicit writes (\texttt{cache[key] = value}) are also accepted
    and bypass the function call.

  \item[Key-range tracking.]
    For integer keys, \texttt{smallest\_key} and \texttt{largest\_key}
    are maintained without iterating the dictionary.
    The \texttt{largest\_key} sentinel drives the tail-convergent iterative
    fill: the recurrence advances from \texttt{largest\_key + 1} to $n$,
    computing only the truly missing entries in $O(\text{gap})$ time and
    $O(1)$ call-stack depth regardless of $n$.

  \item[Usage statistics.]
    \texttt{call\_count} counts cache misses (calls to \texttt{func}),
    \texttt{read\_count} counts cache hits (stored-result retrievals), and
    \texttt{write\_count} counts all writes (including those from seeding
    and intermediate fills).

  \item[Capacity management.]
    An optional \texttt{maxsize} parameter caps the number of stored entries.
    When a new write would exceed the limit, \texttt{prune(n, order)} is
    called automatically using a configurable \texttt{prune\_key} callable.
    The default policy retains the \texttt{maxsize\,-\,1} largest-keyed
    entries, keeping the recurrence frontier intact.  \texttt{reset()} clears
    the cache and all counters; \texttt{clear()} clears entries only, preserving
    \texttt{call\_count} and \texttt{read\_count}.

  \item[Construction-time seeding.]
    The \texttt{seed} constructor parameter pre-populates the cache from a
    plain dict without incrementing any statistics counters.  This is the
    mechanism used by \texttt{CachedGenerator.advance} and
    \texttt{CachedGenerator.insert} to carry relevant entries across
    cache-shifting operations.

  \item[Blocked mutation methods.]
    The mutating \texttt{UserDict} methods (\texttt{pop}, \texttt{popitem},
    \texttt{\_\_delitem\_\_}, \texttt{update}, \texttt{setdefault},
    \texttt{fromkeys}, \texttt{\_\_or\_\_}) are all disabled and raise
    \texttt{NotImplementedError}.  Use \texttt{prune} or \texttt{reset} to
    reduce cache size.
\end{description}

\subsection{Generator-Level Caching: \texttt{CachedGenerator}}

\texttt{CachedGenerator} is a subclass of \texttt{Generator} that internally
holds a \texttt{BaseCache} bound to the wrapped callable.  Every call
\texttt{cgen(n)} is routed through the cache: on a miss the underlying function
is invoked and the result stored; on a hit the stored result is returned directly.

\textbf{When to use each generator type.}
\begin{enumerate}
  \item \texttt{Generator}: the default.  Wrap any callable $g: \N \to \Zp$
        that is cheap to evaluate (e.g.\ a closed-form formula).
  \item \texttt{CachedGenerator}: use when $g(n)$ is expensive \emph{and}
        catena's cache introspection is needed (\texttt{call\_count},
        \texttt{read\_count}, \texttt{stats}, \texttt{BaseCache}
        bookkeeping).  Also the right choice when one expensive generator
        must be shared across several SCF instances so that all instances
        draw from the same memoised pool.
  \item Sharing via \texttt{Generator} wrapping an \texttt{@lru\_cache}-decorated
        function achieves the same functional sharing without catena's
        introspection layer.  Use this pattern for recursive functions
        (Fibonacci, etc.): apply \texttt{@functools.lru\_cache} to the function
        itself \emph{before} passing it to \texttt{Generator} or
        \texttt{CachedGenerator}, because \texttt{CachedGenerator} only
        intercepts top-level calls and does not propagate into self-recursive
        calls inside the function body.
  \item \texttt{FiniteGenerator}: use for fixed finite sequences; no caching
        is needed because lookup is $O(1)$ by array index.
\end{enumerate}

\begin{catenadef}[Caching Patterns]
\begin{lstlisting}[style=python]
from catena.generators import Generator, CachedGenerator
from catena import SimpleContinuedFraction
import functools, math

# --- Pattern 1: expensive non-recursive generator ---
# g(n) = number of divisors of (n+2) not exceeding sqrt(n+2)
def small_divisors(n):
    m = n + 2
    return sum(1 for d in range(1, math.isqrt(m) + 1) if m % d == 0)

cgen = CachedGenerator(small_divisors)
scf = SimpleContinuedFraction(cgen, integer_part=0)
scf.convergent(20)       # fills cache entries 0..20 as a side-effect
cgen.call_count          # 21  (one miss per index)
cgen.read_count          # 0   (no hits yet from this SCF)

# --- Pattern 2: sharing one cache across multiple SCFs ---
scf2 = SimpleContinuedFraction(cgen, integer_part=1)
scf2.convergent(10)      # cgen(0)..cgen(10) are already cached
cgen.read_count          # 11  (all hits: entries 0..10 were populated by scf)

# --- Pattern 3: recursive generator - pre-memo with lru_cache ---
@functools.lru_cache(maxsize=None)
def fib(n):
    return fib(n - 1) + fib(n - 2) if n > 1 else n + 1

# Wrap in CachedGenerator to add BaseCache introspection on top
cgen_fib = CachedGenerator(fib)     # fib's own lru_cache handles recursion
cgen_fib(10)                        # fib(10) was already cached by lru_cache
cgen_fib.call_count                 # 1
cgen_fib.cache.reset()              # clears BaseCache; lru_cache unchanged
\end{lstlisting}
\end{catenadef}

\subsection{Shared Caches and the Tail}

All \texttt{tail()}, \texttt{shift()}, and \texttt{segment(size)} operations
create a new SCF instance via \texttt{\_from\_shared}, which copies the
\emph{reference} to \texttt{\_tail\_cache} rather than duplicating the cache
object.  Consequently both the original and the derived SCF read from and write
to the same \texttt{BaseCache}; a tail-convergent computed by one is
immediately available to the other at zero cost.

This is the mechanism that makes the following hold:

\begin{catenadef}[Shared tail-convergent cache]
\begin{lstlisting}[style=python]
from catena import SimpleContinuedFraction

phi = SimpleContinuedFraction(lambda n: 1, integer_part=1)
phi.tail_convergent(9)             # populates cache entries 0..9

tail = phi.tail()                  # new instance, same _tail_cache
tail.cache_handler is phi.cache_handler   # True
tail.tail_convergent(9)            # O(1) cache hit -- no recomputation
phi.cache_handler.read_count       # incremented by the hit above
\end{lstlisting}
\end{catenadef}